\documentclass[11pt, a4paper]{article}

%% ─── Codificación e idioma ───────────────────────────────────────────────────
\usepackage[utf8]{inputenc}
\usepackage[T1]{fontenc}
\usepackage[spanish, es-nodecimaldot, es-noshorthands]{babel}

%% ─── Tipografía ──────────────────────────────────────────────────────────────
\usepackage{lmodern}
\usepackage{microtype}

%% ─── Geometría de página ─────────────────────────────────────────────────────
\usepackage[top=2.5cm, bottom=2.5cm, left=3cm, right=3cm]{geometry}

%% ─── Matemáticas ─────────────────────────────────────────────────────────────
\usepackage{amsmath}
\usepackage{amssymb}
\usepackage{mathtools}

%% ─── Figuras y tablas ────────────────────────────────────────────────────────
\usepackage{graphicx}
\usepackage{booktabs}
\usepackage{array}
\usepackage{caption}
\usepackage{subcaption}
\usepackage{xcolor}
\usepackage{float}

%% ─── Listas ──────────────────────────────────────────────────────────────────
\usepackage{enumitem}
\setlist[itemize]{noitemsep, topsep=4pt}
\setlist[enumerate]{noitemsep, topsep=4pt}

%% ─── Hipervínculos ───────────────────────────────────────────────────────────
\usepackage[hidelinks, colorlinks=true, linkcolor=blue!70!black,
            citecolor=blue!70!black, urlcolor=blue!70!black]{hyperref}

%% ─── Colores auxiliares ──────────────────────────────────────────────────────
\definecolor{gris}{gray}{0.92}

%% ─── Captions ────────────────────────────────────────────────────────────────
\captionsetup{font=small, labelfont=bf, labelsep=period}

%% ─── Comandos de conveniencia ────────────────────────────────────────────────
\newcommand{\Rx}{\mathrm{Rx}}
\newcommand{\Tx}{\mathrm{Tx}}
\newcommand{\rRx}{\mathbf{r}_{\Rx}}
\newcommand{\dxy}{\delta_{xy}}
\newcommand{\dz}{\delta_z}
\newcommand{\gdeg}{\ensuremath{^{\circ}}}

%% ─── Metadatos ───────────────────────────────────────────────────────────────
\title{\textbf{Análisis de Funciones de Costo para el Optimizador\\
       de Energía/Resolución en SAR Biestático Helicoidal}}
\author{Nota técnica — doctorado SAR}
\date{\today}

% =============================================================================
\begin{document}
% =============================================================================

\maketitle
\thispagestyle{empty}

\noindent\textbf{Contexto:} SAR biestático helicoidal para detección subsuperficial
— optimización de la posición del receptor (Rx) en el plano horizontal.
\vspace{6pt}

\noindent\rule{\linewidth}{0.4pt}

% =============================================================================
\section{Introducción y motivación}
% =============================================================================

El sistema SAR biestático helicoidal opera con el transmisor (Tx) en una trayectoria
circular fija y el receptor (Rx) en una posición configurable. La posición del Rx
determina simultáneamente dos objetivos de calidad de imagen que están en conflicto:

\begin{itemize}
  \item \textbf{Energía recibida} $P_r$: máxima en el ángulo de Brewster
    $\theta_B = \arctan(n_2/n_1)$, donde la transmitancia TM de la interfaz
    suelo–aire es máxima (reflexión nula).

  \item \textbf{Resolución espacial} $(\dxy, \dz)$: controlada por la geometría
    biestática — la resolución lateral mejora cuando el ángulo azimutal entre Tx y Rx
    (medido desde el blanco) $\Delta\phi$ se aleja de cero, pero el ángulo de Brewster
    ocurre a una geometría que no coincide necesariamente con la resolución óptima.
\end{itemize}

La función de costo escalariza este problema biobjetivo en un único escalar a minimizar.
La elección de la función de costo no es neutral: afecta la forma del paisaje de
optimización, la unicidad del mínimo, la sensibilidad al parámetro de compromiso
$\alpha$, y el condicionamiento del algoritmo de búsqueda.

Este documento propone y analiza cinco funciones de costo candidatas con fundamento
en la teoría de optimización multiobjetivo, evaluadas mediante simulación analítica
en una geometría representativa del sistema SAR.

% =============================================================================
\section{Definición formal del problema}
% =============================================================================

\subsection{Espacio de búsqueda y parámetros del sistema}

Sea $\rRx = (x_{\Rx}, y_{\Rx}, z_{\Rx})$ la posición del receptor, con
$z_{\Rx} = 100\,\mathrm{m}$ fijo y $(x_{\Rx}, y_{\Rx}) \in [-270,\,270]^2\,\mathrm{m}$.
Los parámetros del sistema de referencia se listan en la Tabla~\ref{tab:params}.

\begin{table}[H]
  \centering
  \caption{Parámetros del sistema SAR de referencia.}
  \label{tab:params}

  \begin{tabular}{@{}lll@{}}
    \toprule
    \textbf{Parámetro} & \textbf{Valor} & \textbf{Descripción} \\
    \midrule
    $f_0$        & 10\,GHz  & Frecuencia portadora \\
    $B$          & 50\,MHz  & Ancho de banda de barrido \\
    $n_1$        & 1{,}0    & Índice de refracción del aire \\
    $n_2$        & 2{,}0    & Índice del suelo ($\varepsilon_r = 4$) \\
    $\theta_B$   & 63{,}43\gdeg & Ángulo de Brewster TM \\
    $z_{\Tx}$   & 100\,m   & Altitud del Tx \\
    $\rho_{\Tx}$& 160\,m   & Distancia radial del Tx al eje \\
    $B_\mathrm{hel}$ & 47{,}17\,m & Apertura de la hélice \\
    $z_P$        & 5\,m     & Profundidad del blanco \\
    \bottomrule
  \end{tabular}
\end{table}

\subsection{Los dos objetivos elementales}

\paragraph{Objetivo de potencia.}
La potencia recibida en $\rRx$ sigue la ecuación del radar biestático con
refracción TM:

\begin{equation}
  P_r(\rRx) = \frac{P_t \, G_t \, G_r \, \sigma \, T_1 \, T_2(\rRx) \, \lambda_0^2}
                   {(4\pi)^3 \, R_T^2 \, R_R^2(\rRx)}
  \label{eq:Pr}
\end{equation}

donde $T_1$ es la transmitancia TM en el trayecto Tx$\to$blanco (constante para Tx fijo),
$T_2(\rRx)$ es la transmitancia en el trayecto Rx$\to$blanco, y $R_R$ es la longitud
del camino refractado desde Rx hasta el blanco. Ambas cantidades se calculan resolviendo
la condición de Fermat por bisección vectorizada sobre el parámetro de rayo
$p = n_1\sin\theta_{\mathrm{inc}}$.

\paragraph{Objetivo de resolución.}
Las resoluciones derivadas del análisis $k$-espacio del sistema biestático son:

\begin{align}
  \dz &= \frac{c}{2\,W_z}, \qquad
  W_z = n_2 B \cos\theta_{t,0}
        + \frac{c \, B_\mathrm{hel} \sin\psi_0 \cos\psi_0}
               {\lambda_0 \, R_0 \, n_2 \cos\theta_{t,0}}
  \label{eq:dz} \\[6pt]
  \dxy &= \frac{0{,}60\,\lambda_0}{\pi \sin\psi_0 \cdot |\cos(\Delta\phi/2)|}
  \label{eq:dxy}
\end{align}

donde $\psi_0$ es el ángulo de elevación del Rx,
$\theta_{t,0} = \arcsin(\sin\psi_0/n_2)$ el ángulo refractado en el suelo,
$R_0$ la distancia oblicua Rx–superficie, y $\Delta\phi$ el ángulo azimutal
biestático entre Tx y Rx respecto al blanco.

\subsection{Términos objetivos adimensionales}

Para las funciones de costo se definen dos términos escalares adimensionales:

\begin{align}
  f_1(\rRx) &= -\log_{10} P_r(\rRx) \geq 0
  && \text{(pérdida de potencia, a minimizar)}
  \label{eq:f1} \\
  f_2(\rRx) &= \log_{10}\!\bigl[\dxy(\rRx)\cdot\dz(\rRx)\bigr]
  && \text{(producto de resoluciones, a minimizar)}
  \label{eq:f2}
\end{align}

Minimizar $f_1$ equivale a maximizar la potencia recibida; minimizar $f_2$ equivale
a minimizar el volumen de celda de resolución. El conflicto entre ambos objetivos
genera una frontera de Pareto no trivial en el espacio $(f_1, f_2)$.

% =============================================================================
\section{Funciones de costo candidatas}
% =============================================================================

Se proponen cinco funciones de costo con diferentes propiedades teóricas. Cada una
corresponde a una estrategia distinta de escalarización del problema biobjetivo.

% ─────────────────────────────────────────────────────────────────────────────
\subsection{\texorpdfstring{$J_1$}{J1} — Suma ponderada logarítmica (referencia)}
% ─────────────────────────────────────────────────────────────────────────────

\begin{equation}
  \boxed{J_1(\rRx;\,\alpha) = (1-\alpha)\,f_1 + \alpha\,f_2}
  \label{eq:J1}
\end{equation}

\textbf{Interpretación:} escalarización por suma ponderada en escala logarítmica.
El parámetro $\alpha \in [0,1]$ pondera la importancia relativa de la resolución
frente a la potencia.

\textbf{Propiedades:}
\begin{itemize}
  \item Lineal en $f_1$ y $f_2$ $\Rightarrow$ diferenciable en todo el dominio válido.
  \item Ambos términos tienen magnitudes comparables: $f_1 \approx 14$–$18$,
    $f_2 \approx -3$–$-1{,}5$ (la escala logarítmica garantiza conmensurabilidad
    sin normalización explícita).
  \item $\alpha = 0$: solución puramente de potencia (ángulo de Brewster).
  \item $\alpha = 1$: solución puramente de resolución ($\Delta\phi \to \pi$,
    caso degenerado).
  \item Valor recomendado: $\alpha^* \approx 0{,}30$.
  \item \textbf{Limitación:} la suma ponderada solo puede generar puntos en las partes
    convexas de la frontera de Pareto; las regiones no convexas quedan inaccesibles
    para toda combinación de $\alpha$~\cite{miettinen1998}.
\end{itemize}

% ─────────────────────────────────────────────────────────────────────────────
\subsection{\texorpdfstring{$J_2$}{J2} — Suma ponderada lineal normalizada}
% ─────────────────────────────────────────────────────────────────────────────

\begin{equation}
  J_2(\rRx;\,\alpha) = (1-\alpha)\,\tilde{f}_1 + \alpha\,\tilde{f}_2
  \label{eq:J2}
\end{equation}

donde $\tilde{f}_k = (f_k - f_k^{\min}) / (f_k^{\max} - f_k^{\min})$ es la
normalización min-max de cada objetivo sobre el dominio de búsqueda.

\textbf{Propiedades:}
\begin{itemize}
  \item Ambas funciones normalizadas a $[0,1]$: conmensurabilidad garantizada a priori.
  \item Diferenciable (la normalización es una transformación afín).
  \item \textbf{Limitación crítica:} la normalización min-max depende del dominio de
    búsqueda. Si la cuadrícula cambia, la función de costo cambia aunque la física no
    lo haga, lo que viola el principio de invariancia del optimizador.
  \item En escala lineal, $P_r$ varía varios órdenes de magnitud (~$10^{-17}$ a
    $10^{-14}$\,W); en $\tilde{f}_1$ normalizado, la mayor parte de los puntos se
    acumula cerca de cero con sólo el vecindario de Brewster destacando, produciendo
    regiones de gradiente casi nulo (\emph{plateau regions}).
\end{itemize}

% ─────────────────────────────────────────────────────────────────────────────
\subsection{\texorpdfstring{$J_3$}{J3} — Escalarización de Tchebyshev (minimax ponderado)}
% ─────────────────────────────────────────────────────────────────────────────

\begin{equation}
  J_3(\rRx;\,\alpha) = \max\!\bigl[(1-\alpha)\,\tilde{f}_1,\;\alpha\,\tilde{f}_2\bigr]
  \label{eq:J3}
\end{equation}

con la misma normalización min-max que $J_2$.

\textbf{Interpretación:} minimiza la peor desviación ponderada de los objetivos
normalizados respecto a sus valores mínimos individuales. Es la escalarización de
Tchebyshev clásica~\cite{messac2003}, equivalente a la norma $\ell_\infty$ en el
espacio de objetivos normalizados.

\textbf{Propiedades:}
\begin{itemize}
  \item \textbf{Ventaja clave sobre $J_1$:} puede generar puntos en partes no convexas
    de la frontera de Pareto, gracias a la operación de máximo.
  \item No diferenciable donde $(1-\alpha)\tilde{f}_1 = \alpha\tilde{f}_2$: el paisaje
    de optimización tiene aristas, lo que complica métodos basados en gradiente.
  \item La función de Tchebyshev con parámetro de utopía cero garantiza que cada
    minimizador sea eficiente en el sentido de Pareto, independientemente de la
    forma de la frontera.
  \item Variante aumentada: $J_3^{+} = J_3 + \rho\,(\tilde{f}_1 + \tilde{f}_2)$
    con $\rho \ll 1$ es diferenciable y preserva la cobertura completa de
    Pareto~\cite{messac2003}.
\end{itemize}

% ─────────────────────────────────────────────────────────────────────────────
\subsection{\texorpdfstring{$J_4$}{J4} — Distancia euclídea al punto ideal (sin~$\alpha$)}
% ─────────────────────────────────────────────────────────────────────────────

\begin{equation}
  J_4(\rRx) = \sqrt{\tilde{f}_1^{\,2} + \tilde{f}_2^{\,2}}
  \label{eq:J4}
\end{equation}

\textbf{Interpretación:} minimiza la distancia euclídea al origen del espacio de
objetivos normalizados, que corresponde al punto ideal (utopía). No requiere parámetro
de compromiso.

\textbf{Propiedades:}
\begin{itemize}
  \item Genera un único punto del frente de Pareto: el más equitativo en el espacio
    normalizado.
  \item Diferenciable en casi toda parte (excepto en el origen, no alcanzable).
  \item Útil como función de referencia cuando no se dispone de información sobre la
    preferencia relativa entre objetivos.
  \item \textbf{Limitación:} la elección del punto de referencia (origen vs.\ punto
    utopía vs.\ punto anti-utopía) cambia sustancialmente el minimizador, requiriendo
    calibración adicional.
\end{itemize}

% ─────────────────────────────────────────────────────────────────────────────
\subsection{\texorpdfstring{$J_5$}{J5} — Logarítmica con penalización hacia $\theta_B$}
% ─────────────────────────────────────────────────────────────────────────────

\begin{equation}
  J_5(\rRx;\,\alpha,\,\gamma)
  = (1-\alpha)\,f_1 + \alpha\,f_2
    + \gamma\,\bigl(\theta_R - \theta_B\bigr)^2
  \label{eq:J5}
\end{equation}

donde $\theta_R(\rRx)$ es el ángulo de incidencia del rayo Rx$\to$blanco en la
interfaz aire–suelo, y $\theta_B = \arctan(n_2/n_1)$ el ángulo de Brewster.

\textbf{Interpretación:} aumenta $J_1$ con un término de penalización cuadrático que
acerca la solución óptima al ángulo de Brewster. Actúa como restricción blanda:
para $\gamma \to \infty$, la solución converge a la curva $\theta_R = \theta_B$.

\textbf{Propiedades:}
\begin{itemize}
  \item Diferenciable: combina la suavidad de $J_1$ con la penalización cuadrática.
  \item Para $\gamma = 0$ se recupera exactamente $J_1$.
  \item Introduce un tercer hiperparámetro $\gamma \geq 0$, lo que complica la
    calibración respecto a $J_1$.
  \item La penalización es físicamente motivada, pero forzar $\theta_R \to \theta_B$
    puede ser redundante: $J_1$ ya captura el efecto del ángulo de Brewster a través
    de $f_1 = -\log_{10}P_r$, dado que $T_2(\theta_R)$ presenta un máximo agudo
    en $\theta_B$.
\end{itemize}

% =============================================================================
\section{Análisis comparativo}
% =============================================================================

\subsection{Paisajes de optimización}

La Figura~\ref{fig:landscapes} muestra los paisajes de los objetivos elementales
y de las cuatro funciones de costo en el plano $(x_{\Rx}, y_{\Rx})$ para
$\alpha = 0{,}30$ y $\gamma = 5\,\mathrm{rad}^{-2}$.

\begin{figure}[H]
  \centering
  \includegraphics[width=\linewidth]{../utils/fig_cf01_landscapes.pdf}
  \caption{Paisajes de los objetivos elementales ($f_1$, $f_2$) y de las cuatro
    funciones de costo ($J_1$, $J_2$, $J_3$, $J_5$) para $\alpha=0{,}30$.
    Estrella dorada: posición del Tx.
    Cruz blanca: proyección del blanco en superficie.
    Círculo blanco: posición óptima del Rx (mínimo de cada función).}
  \label{fig:landscapes}
\end{figure}

\textbf{Observaciones:}
\begin{enumerate}
  \item El paisaje de $f_1$ (potencia) exhibe un anillo de alta transmitancia alrededor
    del blanco, correspondiente al ángulo de Brewster $\theta_B = 63{,}4\gdeg$. El pico
    de transmitancia TM es angosto (${\sim}5\gdeg$) y concentra el gradiente de $P_r$ en
    esa región angular.

  \item El paisaje de $f_2$ (resolución) muestra que la resolución lateral $\dxy$
    mejora a medida que el ángulo azimutal $\Delta\phi$ entre Tx y Rx se aleja de
    cero, con mínimo en $\Delta\phi = 180\gdeg$ (Rx diametralmente opuesto al Tx).
    La resolución axial $\dz$ depende principalmente de la distancia radial $\rho_{\Rx}$.

  \item Los mínimos de $J_1$ y $J_5$ son similares en ubicación, con $J_5$ desplazando
    ligeramente el óptimo hacia la curva de Brewster. El paisaje de $J_3$ (Tchebyshev)
    muestra aristas características de la operación $\max$, con mínimo en una región
    más compacta.

  \item $J_2$ (lineal normalizada) exhibe un paisaje más plano fuera de la vecindad
    del ángulo de Brewster, lo que dificulta la convergencia de métodos de gradiente.
\end{enumerate}

\subsection{Fronteras de Pareto}

La Figura~\ref{fig:pareto} muestra las fronteras de Pareto trazadas por cada función
al variar $\alpha \in [0,1]$ con paso fino.

\begin{figure}[H]
  \centering
  \includegraphics[width=0.8\linewidth]{../utils/fig_cf02_pareto_fronts.pdf}
  \caption{Fronteras de Pareto en el espacio de objetivos $(f_1, f_2)$.
    Cada punto corresponde al óptimo de la función para un valor de $\alpha$.
    El punto ideal (estrella negra) es el mínimo simultáneo teórico de ambos objetivos.
    El rombo ($J_4$) es el punto de mínima distancia euclídea al punto ideal.}
  \label{fig:pareto}
\end{figure}

\textbf{Observaciones:}
\begin{enumerate}
  \item $J_1$ genera una curva de Pareto relativamente suave. Pequeñas variaciones de
    $\alpha$ en $[0{,}2,\,0{,}4]$ producen desplazamientos significativos en la posición
    óptima del Rx, indicando alta sensibilidad en esa región.

  \item $J_3$ (Tchebyshev) traza una curva de Pareto diferente: cubre partes del frente
    que $J_1$ no puede alcanzar con ningún $\alpha$ (regiones no convexas), confirmando
    la ventaja teórica de la escalarización de Tchebyshev~\cite{das1997}.

  \item $J_2$ traza una curva cuya posición relativa depende del dominio de normalización,
    haciendo la comparación directa con $J_1$ y $J_3$ no independiente del dominio.

  \item $J_4$ (distancia euclídea) produce un único punto ubicado cerca del codo de la
    curva de Pareto de $J_1$, coherente con el objetivo de mínima distancia al punto ideal.
\end{enumerate}

\subsection{Sensibilidad al parámetro \texorpdfstring{$\alpha$}{alpha}}

La Figura~\ref{fig:alpha} muestra la posición óptima del Rx (distancia radial y azimut)
como función de $\alpha$.

\begin{figure}[H]
  \centering
  \includegraphics[width=\linewidth]{../utils/fig_cf03_alpha_sensitivity.pdf}
  \caption{(a) Distancia radial $\rho_{\Rx}^*$ del Rx óptimo al blanco vs.\ $\alpha$.
    (b) Azimut $\phi_{\Rx}^*$ del Rx óptimo vs.\ $\alpha$.
    Se comparan $J_1$, $J_2$ y $J_3$.}
  \label{fig:alpha}
\end{figure}

\textbf{Observaciones:}
\begin{enumerate}
  \item Para $J_1$, la distancia radial óptima $\rho_{\Rx}^*$ varía de forma monótona
    suave con $\alpha$: para $\alpha \to 0$ el Rx tiende hacia el anillo de Brewster
    ($\rho \propto z_{\Rx}\tan\theta_B \approx 200\,\mathrm{m}$); para $\alpha \to 1$
    el Rx se desplaza radialmente en busca de mejor resolución axial.

  \item El azimut $\phi_{\Rx}^*$ de $J_1$ varía también de forma continua: para $\alpha$
    pequeño, el Rx se ubica en el lado del anillo de Brewster opuesto al Tx; para
    $\alpha$ grande, la resolución lateral domina y el Rx tiende a $\Delta\phi \to 180\gdeg$.

  \item $J_3$ exhibe saltos abruptos (discontinuidades) en $\rho_{\Rx}^*$ y $\phi_{\Rx}^*$
    a ciertos valores de $\alpha$. Esto es consecuencia de la operación $\max$: al cambiar
    cuál de las dos ramas domina, el argmín puede saltar a otra región del dominio.
    Esta característica es una propiedad inherente de la escalarización de Tchebyshev,
    no un artefacto numérico.

  \item La región $\alpha \in [0{,}25,\,0{,}35]$ es la de mayor variabilidad para $J_1$,
    correspondiendo a la transición entre el régimen dominado por Brewster y el dominado
    por la geometría biestática.
\end{enumerate}

\subsection{Condicionamiento del gradiente}

La Figura~\ref{fig:grad} muestra la norma del gradiente $\|\nabla J\|$ a lo largo
de un perfil horizontal y el mapa de $\|\nabla J_1\|$.

\begin{figure}[H]
  \centering
  \includegraphics[width=\linewidth]{../utils/fig_cf04_gradient_norm.pdf}
  \caption{(a) $\|\nabla J\|$ normalizado a lo largo del perfil $y_{\Rx}=0$ para
    $\alpha=0{,}30$. (b) Mapa de $\|\nabla J_1\|$ con isolíneas de $J_1$ superpuestas.}
  \label{fig:grad}
\end{figure}

\textbf{Observaciones:}
\begin{enumerate}
  \item $J_1$ tiene gradiente significativo en toda la región de interés, incluyendo
    zonas alejadas del mínimo. Esto favorece la convergencia de métodos de gradiente
    descendente desde posiciones iniciales diversas.

  \item $J_2$ exhibe regiones amplias con gradiente casi nulo, especialmente lejos del
    ángulo de Brewster (\emph{plateau regions}~\cite{nocedal2006}), lo que dificulta
    la convergencia de métodos locales de primera derivada.

  \item $J_3$ presenta gradiente alto pero con discontinuidades locales en las aristas
    de la operación $\max$. El panel~(b) muestra cómo los gradientes más altos de
    $J_1$ se concentran en el anillo de Brewster, coherente con la fuerte variación
    de $T_2(\theta_R)$ en esa región angular.

  \item El gradiente de $J_1$ es más suave y consistente. Su número de condición
    $\kappa = \|\nabla^2 J\|_F \,/\, \lambda_{\min}(\nabla^2 J)$ es más favorable
    que el de $J_2$, lo que se traduce en mejor tasa de convergencia para métodos
    cuasi-Newton o Newton-Raphson.
\end{enumerate}

\subsection{Efecto de la penalización de Brewster (\texorpdfstring{$J_5$}{J5})}

La Figura~\ref{fig:brewster} analiza el efecto del parámetro de penalización $\gamma$
en la función $J_5$.

\begin{figure}[H]
  \centering
  \includegraphics[width=\linewidth]{../utils/fig_cf05_brewster_penalty.pdf}
  \caption{(a) Ángulo de incidencia $\theta_R^*$ en el Rx óptimo vs.\ $\gamma$,
    con la asíntota $\theta_B$ (línea discontinua). (b) Pérdida de potencia $f_1$
    y logaritmo del producto de resolución $f_2$ en el Rx óptimo vs.\ $\gamma$.}
  \label{fig:brewster}
\end{figure}

\textbf{Observaciones:}
\begin{enumerate}
  \item Para $\gamma = 0$, $J_5 \equiv J_1$. Al aumentar $\gamma$, el Rx óptimo se
    desplaza hacia la posición donde $\theta_R = \theta_B$. La convergencia es rápida
    para $\gamma < 5\,\mathrm{rad}^{-2}$ y satura para $\gamma > 10\,\mathrm{rad}^{-2}$.

  \item Panel~(b): el aumento de $\gamma$ reduce $f_1$ (Brewster maximiza $P_r$) pero
    aumenta $f_2$ (la posición de Brewster no es la de resolución óptima), confirmando
    el sesgo físicamente justificado pero costoso en resolución.

  \item El rango $\gamma \in [1,\,5]\,\mathrm{rad}^{-2}$ ofrece un compromiso razonable:
    el Rx converge a ${\sim}85\%$ del camino hacia $\theta_B$ con degradación de
    resolución inferior al $15\%$ respecto al caso $\gamma = 0$.
\end{enumerate}

% =============================================================================
\section{Tabla comparativa}
% =============================================================================

\begin{table}[H]
  \centering
  \caption{Comparación de las cinco funciones de costo candidatas.}
  \label{tab:comparison}
  \renewcommand{\arraystretch}{1.35}

  \begin{tabular}{@{} p{4.8cm} c c c c c @{}}
    \toprule
    \textbf{Propiedad}
      & $J_1$ & $J_2$ & $J_3$ & $J_4$ & $J_5$ \\
    \midrule
    Escala
      & Log. & Lineal & Lineal & Lineal & Log. \\
    Diferenciable
      & Sí & Sí & No & Sí & Sí \\
    Parámetros
      & $\alpha$ & $\alpha$ & $\alpha$ & — & $\alpha,\gamma$ \\
    Cubre frente no convexo
      & No & No & \textbf{Sí} & Parcial & No \\
    Normalización dependiente del dominio
      & No & \textbf{Sí} & \textbf{Sí} & \textbf{Sí} & No \\
    Condicionamiento numérico
      & \textbf{Bueno} & Pobre & Moderado & Moderado & Bueno \\
    Paisaje unimodal
      & \textbf{Sí} & Sí & Sí & Sí & \textbf{Sí} \\
    Interpretación física directa
      & \textbf{Sí} & Moderada & Moderada & Baja & Sí \\
    Calibración requerida
      & $\alpha$ & $\alpha$+escala & $\alpha$+escala & Escala & $\alpha,\gamma$ \\
    \bottomrule
  \end{tabular}
\end{table}

% =============================================================================
\section{Discusión}
% =============================================================================

\subsection{Ventajas de la formulación logarítmica (\texorpdfstring{$J_1$}{J1})}

La función de costo actual $J_1$ exhibe una propiedad notable que no es inmediatamente
obvia: la escala logarítmica actúa como normalizador intrínseco. Dado que $P_r$ varía
en varios órdenes de magnitud a lo largo del dominio de búsqueda, su representación
logarítmica comprime la variabilidad de modo que $f_1 \in [14,\,18]$ y
$f_2 \in [-3,\,-1{,}5]$ resultan magnitudes comparables sin normalización explícita.
Esto elimina la dependencia del dominio de búsqueda que afecta a $J_2$, $J_3$ y $J_4$.

Adicionalmente, la función logarítmica es cóncava: garantiza que los puntos de mayor
potencia (ángulo de Brewster) estén suficientemente separados en el espacio de costo
para generar gradiente apreciable, mientras que variaciones pequeñas de $P_r$ lejos
de $\theta_B$ contribuyen modestamente a $f_1$.

\subsection{Cuándo conviene explorar \texorpdfstring{$J_3$}{J3} (Tchebyshev)}

Si el análisis de la frontera de Pareto revela que la geometría del sistema produce
una frontera no convexa — lo cual puede ocurrir cuando la transmitancia TM exhibe
múltiples máximos locales (e.g., en estructuras multicapa con capas resonantes de
media longitud de onda) — entonces $J_1$ no puede alcanzar todos los puntos eficientes
en el sentido de Pareto. En esos escenarios, la escalarización $J_3$ es preferible,
aunque a costa de un paisaje no diferenciable que requiere métodos de búsqueda global
(algoritmos genéticos, \emph{simulated annealing}) o métodos de subgradiente.

Para el sistema de dos capas (aire–suelo) aquí analizado, la frontera de Pareto es
convexa o suavemente no convexa, por lo que $J_1$ es suficiente.

\subsection{Uso de \texorpdfstring{$J_5$}{J5} como restricción blanda}

La función $J_5$ es conceptualmente interesante como herramienta de diseño: permite
explorar el compromiso entre la importancia de operar en el ángulo de Brewster y la
pérdida de resolución que esto conlleva. Sin embargo, para el optimizador de producción
introduce una complejidad adicional (selección de $\gamma$) sin ventaja clara sobre la
elección directa de $\alpha$ pequeño, que ya privilegia $f_1 = -\log_{10}P_r$.

Una alternativa más rigurosa al uso de $J_5$ es la formulación $\varepsilon$-constrained:

\begin{equation}
  \min_{\rRx} \; f_2(\rRx) \quad
  \text{sujeto a} \quad f_1(\rRx) \leq f_1^{\max}
  \label{eq:eps_constrained}
\end{equation}

donde $f_1^{\max}$ es un umbral de pérdida de potencia aceptable. Esta formulación
garantiza que la potencia mínima requerida sea estrictamente satisfecha en lugar de
ser penalizada suavemente.

% =============================================================================
\section{Conclusiones y recomendación}
% =============================================================================

Con base en el análisis teórico y las simulaciones presentadas, se establecen las
siguientes conclusiones:

\begin{enumerate}
  \item \textbf{La función logarítmica $J_1$ es la opción óptima} para el sistema SAR
    biestático helicoidal de dos capas (aire–suelo). Su conmensurabilidad intrínseca,
    buena diferenciabilidad y paisaje unimodal la hacen compatible con métodos de
    optimización local eficientes. El valor recomendado $\alpha^* = 0{,}30$ equilibra
    potencia y resolución en la región de mayor sensibilidad de la frontera de Pareto.

  \item \textbf{La escalarización de Tchebyshev $J_3$} debe considerarse si el sistema
    se extiende a medios multicapa con capas resonantes, donde la frontera de Pareto
    puede ser no convexa. Se recomienda entonces la variante aumentada con
    $\rho = 10^{-3}$ para recuperar diferenciabilidad.

  \item \textbf{La función lineal normalizada $J_2$} no se recomienda para uso en
    producción por su dependencia del dominio de búsqueda y su condicionamiento
    numérico pobre. Puede ser útil como verificación cualitativa de la convexidad del
    frente de Pareto.

  \item \textbf{La penalización de Brewster $J_5$} con $\gamma \in [1,5]\,\mathrm{rad}^{-2}$
    ofrece un mecanismo físicamente interpretable para priorizar la transmitancia TM,
    pero se recomienda sustituirla por la formulación $\varepsilon$-constrained~%
    \eqref{eq:eps_constrained} cuando la pérdida de potencia mínima sea un requisito
    de diseño estricto.

  \item Para la extensión a sistemas multicapa ($N > 2$ capas), la formulación de $J_1$
    se generaliza directamente: basta reemplazar $T_2(\theta_R)$ con la transmitancia
    total del \emph{stack} ABCD evaluada en el ángulo refractado desde Rx. La función
    matemática no cambia; sólo cambia el modelo de propagación subyacente.
\end{enumerate}

% =============================================================================
%  BIBLIOGRAFÍA
% =============================================================================
\begin{thebibliography}{9}

\bibitem{miettinen1998}
K.~Miettinen,
\textit{Nonlinear Multiobjective Optimization}.
Springer, 1998, cap.~3.

\bibitem{messac2003}
A.~Messac, A.~Ismail-Yahaya, and C.~A.~Mattson,
``The normalized normal constraint method for generating the Pareto frontier,''
\textit{Struct.\ Multidisc.\ Optim.}, vol.~25, no.~2, pp.~86--98, 2003.

\bibitem{nocedal2006}
J.~Nocedal and S.~J.~Wright,
\textit{Numerical Optimization}, 2nd~ed.
Springer, 2006, pp.~30--35.

\bibitem{das1997}
I.~Das and J.~E.~Dennis,
``A closer look at drawbacks of minimizing weighted sums of objectives
for Pareto set generation in multicriteria optimization problems,''
\textit{Struct.\ Optim.}, vol.~14, pp.~63--69, 1997.

\bibitem{moreira2013}
A.~Moreira \textit{et al.},
``A tutorial on synthetic aperture radar,''
\textit{IEEE Geosci.\ Remote Sens.\ Mag.}, vol.~1, no.~1, pp.~6--43, Mar.~2013.

\end{thebibliography}

\bigskip
\noindent\rule{\linewidth}{0.4pt}

\noindent\small
\textbf{Script:} \texttt{utils/generate\_figures\_cost\_analysis.py}
\begin{itemize}[noitemsep, topsep=2pt, leftmargin=1.5em]
  \item[] \texttt{utils/fig\_cf01\_landscapes.pdf}
  \item[] \texttt{utils/fig\_cf02\_pareto\_fronts.pdf}
  \item[] \texttt{utils/fig\_cf03\_alpha\_sensitivity.pdf}
  \item[] \texttt{utils/fig\_cf04\_gradient\_norm.pdf}
  \item[] \texttt{utils/fig\_cf05\_brewster\_penalty.pdf}
\end{itemize}

\end{document}
